\chapter{Reduction and boundary fields}
\label{mht:ch:neumann-bv}

Putting the curve connection in normal form also changes the surface
connection. Before imposing the field equations, the transformed surface
field has an error determined by the mixed curvature. On flat fields
this error vanishes, giving the full formal BF reduction. At a fixed
residual matrix, the linearized reduction retains its entire adjoint
complex, including the stabilizer directions that appear when
eigenvalues coincide.

The action retains two kinds of boundary data. At the pole, its Stokes
term pairs the value of a function with its first holomorphic jet.
Along the boundary of the surface, a varying gauge transformation
contributes a cubic term in its Maurer--Cartan form. We derive both
terms after computing the exact quadratic action of the eliminated
surface fields.

\section{Curvature and the surface connection}

Keep the sphere, trace, normalized Green operator, and representative
$f_0,\eta_0$ of Chapter~\ref{mht:ch:polar-neumann}.
Write $\mathfrak g=\mathfrak{gl}_N(\mathbb C)$ and let $U$ be an oriented
smooth surface. Surface forms precede curve forms in the graded tensor
product. Decompose a physical mixed connection as
\[
 \mathsf A=b+a,\qquad
 b\in\Omega^1(U)\widehat\otimes E^0\otimes\mathfrak g,\qquad
 a\in C^\infty(U)\widehat\otimes E^1\otimes\mathfrak g.
\]

Work formally in the field amplitude, or in the small Sobolev
neighborhood of Proposition~\ref{mht:prop:analytic-frame}, smoothly
parametrized over $U$. Let $H=H(a)$ and $J=\mathcal J(a)$, so
$\bar\partial H=-aH$, $H(o)=I_N$, and $j(H)=-J$.
Put $b_o=b(o)$.

Transformation by $H$ kills the curve connection and leaves the surface
one-form
\begin{equation}\label{mht:eq:surface-frame}
 B=H^{-1}bH+H^{-1}d_UH.
\end{equation}
To recover it, define the curvature components and their residual value:
\begin{align}
 F_{UU}&=d_Ub+b^2,&
 F_{UC}&=d_Ua-\bar\partial b+[b,a],\label{mht:eq:mixed-curvature}\\
 C_a&=H^{-1}F_{UC}H,&
 r&=d_UJ+[b_o,J].\notag
\end{align}
Here $\bar\partial b$ differentiates the curve coefficients. Its minus
sign is the tensor sign in the total differential.

\begin{proposition}\label{mht:prop:curvature-reconstruction}
Without imposing flatness,
\begin{equation}\label{mht:eq:curvature-reconstruction}
 B=b_o-GC_a,\qquad
 r=j(GC_a)=\kappa^{-1}\tau_E(C_a),\qquad
 (H^{-1}F_{UU}H)(o)=d_Ub_o+b_o^2.
\end{equation}
All curve operators act coefficientwise on surface forms.
\end{proposition}

\begin{proof}
The transformed curve connection is zero, so covariance of mixed
curvature gives $-\bar\partial B=C_a$. Since $B(o)=b_o$, normalization
of $G$ gives the first identity. Differentiate
\eqref{mht:eq:surface-frame} at $o$, using
$j(H^{-1})=J$, $j(b)=0$, and $j(d_UH)=-d_UJ$.
The result is $j(B)=-r$, which gives the second identity by polar Stokes.
Surface curvature transforms by conjugation, and evaluation at $o$
gives the last identity.
\end{proof}

The frame $H$ need not be Neumann. The admissible transformation is
$g=H\exp(f_0J)$ from Theorem~\ref{mht:thm:neumann-normal-form}.
For $L=\operatorname{ad}_J$, introduce the entire operator
\begin{equation}\label{mht:eq:entire-adjoint-operator}
 T_J=f_0\int_0^1e^{-u f_0L}\,du
       =\sum_{n\geq0}\frac{(-1)^nf_0^{n+1}L^n}{(n+1)!}.
\end{equation}
This formula is defined even when $L$ has a kernel.

\begin{corollary}\label{mht:cor:full-normal-error}
For the convention
$\mathsf A^g=g^{-1}\mathsf A g+g^{-1}(d_U+\bar\partial)g$,
\begin{equation}\label{mht:eq:full-normal-error}
 \mathsf A^g=b_o+\eta_0J+R_a,\qquad
 R_a=-e^{-f_0L}GC_a+T_Jr.
\end{equation}
The surface one-form $R_a$ has value and first holomorphic jet zero at
$o$.
\end{corollary}

\begin{proof}
After transformation by $H$, transform $B$ by $e^{f_0J}$.
Differentiation of the matrix exponential gives
\[
 e^{-f_0J}Be^{f_0J}+e^{-f_0J}d_Ue^{f_0J}-b_o
   =e^{-f_0L}(B-b_o)+T_J(d_UJ+[b_o,J]).
\]
Substitute \eqref{mht:eq:curvature-reconstruction}.
The curve component becomes $\eta_0J$.
Both $GC_a$ and $T_J$ vanish at $o$, and their first jets are $r$ and
the identity operator. Their contributions to $j(R_a)$ cancel.
\end{proof}

Thus mixed curvature measures exactly the failure of a general field
to be represented by residual BF fields. On flat fields, $C_a=r=0$,
and the normal form is $b_o+\eta_0J$.

\section{Formal fields and gauge homotopies}

To retain ghosts and higher gauge homotopies, let $R$ be a local Artin
commutative differential graded algebra over $\mathbb C$, with nilpotent
maximal ideal $\mathfrak m$. A formal field is a total-degree-one element
of
\[
 \Omega^\bullet(U)\widehat\otimes E\otimes\mathfrak g
                                      \otimes\mathfrak m.
\]
It is a Maurer--Cartan field if $D\mathsf A+\mathsf A^2=0$, where $D$
is the total differential. Tensoring additionally with
$\Omega^\bullet_{\mathrm{poly}}(\Delta^n)$ defines its simplicial space
of fields. The residual complex replaces $E$ by
$A_\kappa=\mathbb C\oplus\mathbb C\varepsilon$.

\begin{theorem}\label{mht:thm:formal-bf-reduction}
The maps
\[
 \mathsf A=\mathbf b+a\longmapsto(\mathbf b(o),\mathcal J(a)),
 \qquad
 (\mathbf b_o,\mathbf J)\longmapsto\mathbf b_o+\mathbf J\eta_0
\]
give mutually inverse equivalences of these formal deformation
problems, up to a natural based Neumann gauge homotopy.
The residual equations are
\begin{equation}\label{mht:eq:residual-mc}
 d\mathbf b_o+\mathbf b_o^2=0,\qquad
 d\mathbf J+[\mathbf b_o,\mathbf J]=0.
\end{equation}
Here $d$ includes the surface, Artin, and simplex differentials.
\end{theorem}

\begin{proof}
The coefficient of a curve one-form has total coefficient degree zero.
Its normalized frame therefore has degree zero, and the Green series
and exponential are defined over $R$. Nilpotence makes the series
finite. The curvature calculation uses only the Leibniz rule and graded
matrix multiplication, so it applies with the full coefficient
differential. On a Maurer--Cartan field, both curvature components
vanish. Equation~\eqref{mht:eq:full-normal-error} then gives the residual
field and \eqref{mht:eq:residual-mc}.

For a residual field, $H=\exp(-f_0\mathbf J)$, so the residual composite
is the identity and the correcting gauge is $g=I_N$.
For a general field, set $g_u=\exp(u\log g)$.
The connection
\[
 \mathsf A^{g_u}+g_u^{-1}\partial_ug_u\,du
\]
is Maurer--Cartan on the parameter interval. It joins $\mathsf A$ to
its residual normal form. Its residual image is constant, since
$g_u(o)=I_N$, $j(g_u)=0$, and its additional $du$ coefficient vanishes
at $o$.

For a Neumann gauge $k$, uniqueness of the frame gives
$H(a^k)=k^{-1}H(a)k(o)$. Hence $\mathcal J$ transforms by conjugation
with $k(o)$, and $\mathbf b(o)$ transforms as a connection.
The formulas commute with coefficient morphisms and simplex
restrictions. They consequently give the stated equivalence with its
higher gauge homotopies.
\end{proof}

The strict inclusion preserves the cyclic pairing. On a closed surface,
or with compact support, substitution in the cyclic action gives
\begin{align}
 S_{\mathrm{BF}}
   &=\kappa\int_U\operatorname{tr}
             \bigl(\mathbf J(d_U\mathbf b+\mathbf b^2)\bigr),
       \label{mht:eq:residual-bv-action}\\
 \omega_{\mathrm{BF}}
   &=\kappa\int_U\operatorname{tr}
                      (\delta\mathbf J\,\delta\mathbf b).
       \label{mht:eq:residual-bv-pairing}
\end{align}
The pairing has degree $-1$. In surface degrees $0,1,2$, write
$\mathbf b=(c,b,J^+)$ and $\mathbf J=(J,b^+,c^+)$.
Their coordinate ghost numbers are respectively $(1,0,-1)$ and
$(0,-1,-2)$. These are all components of the residual shifted complex.

The projection removes contractible directions, so it is not a
symplectomorphism of the full spaces of fields before the equations of
motion are imposed. The equivalence concerns the trivial curve-bundle
neighborhood. Its maps preserve restriction and support in $U$, while
the Green operator still uses the whole curve.

\section{The adjoint complex at a residual matrix}

The residual field $J$ may have repeated eigenvalues or be nonsemisimple.
The correct linearized complex keeps both copies of $\mathfrak g$:
\begin{equation}\label{mht:eq:residual-adjoint-complex}
 K_J=\bigl[\mathfrak g\xrightarrow{\operatorname{ad}_J}
                                  \mathfrak g\varepsilon\bigr].
\end{equation}
They lie in degrees zero and one. On the smooth curve complex, the
differential is
\[
 d_Jf=\bar\partial f+\eta_0[J,f],\qquad d_J(E^1)=0.
\]

The inclusion $I:K_J\to E\otimes\mathfrak g$ is
$I_0x=x$, $I_1y=\eta_0y$. Extend the cyclic projections entrywise:
\[
 P_0f=\kappa^{-1}\tau_E(f\eta_0),\qquad
 P_1\eta=\kappa^{-1}\tau_E(\eta).
\]
The factor $\varepsilon$ in the target of $P_1$ is understood.

Conjugation by $e^{f_0\operatorname{ad}_J}$ removes the twist from
$d_J$. It changes the domain condition, so a correction is still
needed. With $L=\operatorname{ad}_J$ and $T_J$ from
\eqref{mht:eq:entire-adjoint-operator}, define
\begin{align}
 p^{\mathrm{fr}}_{J,0}f&=f(o),&
 p^{\mathrm{fr}}_{J,1}\eta&=jG(e^{f_0L}\eta),\notag\\
 h^{\mathrm{fr}}_J\eta
   &=e^{-f_0L}G(e^{f_0L}\eta)-T_Jp^{\mathrm{fr}}_{J,1}\eta,
       &h^{\mathrm{fr}}_J(E^0)&=0.
       \label{mht:eq:background-frame-h}
\end{align}

\begin{theorem}\label{mht:thm:background-contraction}
For every finite matrix $J$, the map
\begin{equation}\label{mht:eq:background-cyclic-h}
 h_J=(1-I_0P_0)h^{\mathrm{fr}}_J
\end{equation}
is a continuous cyclic contraction onto $K_J$:
\[
 PI=1,\qquad d_Jh_J+h_Jd_J=1-IP,\qquad
 h_J^2=h_JI=Ph_J=0.
\]
For matrix-valued curve one-forms $\eta,\xi$,
\begin{equation}\label{mht:eq:background-cyclic-identity}
 \tau_E\operatorname{tr}((h_J\eta)\xi)
       +\tau_E\operatorname{tr}(\eta h_J\xi)=0.
\end{equation}
The maps depend analytically on $J$, including at nonsemisimple
matrices, and define coefficientwise maps for formal matrix parameters.
The cohomology is $\ker\operatorname{ad}_J$ in degree zero and
$\operatorname{coker}\operatorname{ad}_J$ in degree one.
\end{theorem}

\begin{proof}
For $F=e^{f_0L}f$, the Neumann condition $j(f)=0$ becomes
$j(F)=LF(o)$. Direct differentiation on unconstrained smooth functions
gives
\[
 d_J(T_Jy)=\eta_0y,\qquad T_JL=1-e^{-f_0L}.
\]
The two first jets in \eqref{mht:eq:background-frame-h} cancel.
Thus $h^{\mathrm{fr}}_J$ takes values in $E^0$ and vanishes at $o$.
The same identities give
\[
 d_Jh^{\mathrm{fr}}_J\eta
       =\eta-\eta_0p^{\mathrm{fr}}_{J,1}\eta,\qquad
 h^{\mathrm{fr}}_Jd_Jf=f-f(o).
\]
For the latter equality use
$G(e^{f_0L}d_Jf)=e^{f_0L}f-f(o)$.

The normalized primitive of $e^{f_0L}\eta_0y$ is
$\int_0^{f_0}e^{uL}y\,du$. It proves
$p^{\mathrm{fr}}_JI=1$ and $h^{\mathrm{fr}}_JI=0$.
The adjoint projection is a chain map, since
$P_1d_Jf=[J,P_0f]$. Hence $\Pi=1-IP$ commutes with $d_J$, and
$h_J=\Pi h^{\mathrm{fr}}_J\Pi$. Therefore
\[
 d_Jh_J+h_Jd_J
   =\Pi(1-Ip^{\mathrm{fr}}_J)\Pi=\Pi.
\]
The side conditions follow from $\Pi I=P\Pi=0$,
$h^{\mathrm{fr}}_JI=0$, and degree.

For cyclicity, put $v=h_J\eta$ and $w=h_J\xi$.
Both have zero $P_0$. Adjointness of $P$ and $I$, followed by polar
integration by parts, gives
\[
 \tau_E\operatorname{tr}(v\xi)
 =\tau_E\operatorname{tr}(v\,d_Jw)
 =-\tau_E\operatorname{tr}((d_Jv)w)
 =-\tau_E\operatorname{tr}(\eta w).
\]
The commutator contribution has zero matrix trace, and the degree-zero
products remain in the Neumann domain.

Finally, multiplication by the smooth matrix exponentials, $G$, and the
finite-rank projections is continuous. Parameter derivatives are
bounded on compact sets of matrices in every smooth seminorm.
This proves the analytic assertion and the coefficientwise formal
construction. The contraction identifies cohomology with that of
\eqref{mht:eq:residual-adjoint-complex}.
\end{proof}

In particular, put
\[
 V=\ker P_0\subset E^0\otimes\mathfrak g,\qquad
 W=\ker P_1\subset E^1\otimes\mathfrak g.
\]
Then $d_J:V\to W$ is an isomorphism with inverse $h_J$.
Write $h=h_{J=0}$ for the cyclic homotopy of
Theorem~\ref{mht:thm:cyclic-contraction}.
As a formal series at $J=0$,
\[
 h_J=h(1+\eta_0\operatorname{ad}_Jh)^{-1}.
\]
The exact formula \eqref{mht:eq:background-cyclic-h} does not require
convergence of this particular series.

\subsection{An eigenvalue collision}

The family $J_x=\operatorname{diag}(x,-x)$ shows what would be lost by
dividing by nonzero eigenvalue differences. Over either
$R=\mathbb C[x]$ or $R=\mathbb C[[x]]$, its residual complex has
\begin{equation}\label{mht:eq:collision-torsion}
 H^0(K_{J_x})=R^2,\qquad
 H^1(K_{J_x})=R^2\oplus(R/(x))^2.
\end{equation}
Indeed, the two diagonal directions have zero differential, and the
two off-diagonal arrows are multiplication by $2x$ and $-2x$.
They are injective over $R$ with cokernel $R/(x)$.

At $x=0$, the specialized complex has four-dimensional kernel and
cokernel. Each two-term free resolution of $R/(x)$ contributes a
$\operatorname{Tor}_1^R(R/(x),\mathbb C)$ class in degree zero.
This accounts for the two additional stabilizer directions.
Theorem~\ref{mht:thm:background-contraction} transports the calculation
to the native formal family $(E\otimes\mathfrak g)[[x]]$.

The completion in this last assertion is necessary. For
$f_0=z/(1+z\bar z)$, the differential in the $E_{12}$ direction is
$D_x=\bar\partial+2x\eta_0$ on $E^0[x]$. The cokernel of the residual
inclusion on first cohomology is
\[
 Q=\frac{E^1[x]}{D_xE^0[x]+\mathbb C[x]\eta_0}.
\]

\begin{proposition}\label{mht:prop:polynomial-family-failure}
Inclusion as constant polynomials induces an injective complex linear
map $E^1/\mathbb C\eta_0\longrightarrow Q$.
For $v=z\bar z/(1+z\bar z)$, the classes
$[\bar\partial(v^m)]$, $m\geq1$, are linearly independent in $Q$.
Thus the residual inclusion over $\mathbb C[x]$ has an
infinite-dimensional cokernel on first cohomology of
$(E\otimes\mathfrak{gl}_2)[x]$ with differential $d_{J_x}$.
\end{proposition}

\begin{proof}
Suppose that $\eta\in E^1$, viewed as constant in $x$, has zero class in $Q$.
There are polynomials $c(x)\in\mathbb C[x]$ and $u(x)\in E^0[x]$ with
\[
 \eta=c(x)\eta_0+D_xu(x).
\]
Multiplication by $\eta_0$ on smooth functions is injective because the
form is nonzero on a dense open set.

If $u$ has degree $n$, comparison above degree $n+1$ first gives
$\deg c\leq n+1$. The coefficient of $x^{n+1}$ says that the top
coefficient of $u$ is a scalar constant. Its Dolbeault derivative
vanishes. Descending through powers of $x$ then makes every coefficient
of $u$ constant. The constant term gives $\eta=c(0)\eta_0$,
which also follows directly when $u=0$. This proves the injection.

A finite relation among the displayed classes would therefore give
\[
 \bar\partial\left(\sum_{m=1}^M a_m v^m-cf_0\right)=0.
\]
The function in parentheses is smooth on the sphere, so it is
holomorphic and constant. Its first holomorphic jet at $o$ is $-c$,
so $c=0$, and its value at $o$ makes the constant zero.
Since $v$ takes every value in $[0,1)$ on the affine chart,
the polynomial $\sum_{m=1}^M a_m t^m$ vanishes identically.
Every coefficient $a_m$ is zero.
\end{proof}

Thus a formal or analytic parameter family retains the collision through
the full adjoint complex. The exponential contraction is not a map on
the uncompleted polynomial family.

\section{The exact action of the eliminated fields}

On the curve slice $a=\eta_0J$, split the surface field by the cyclic
projection:
\[
 b=b_c+u,\qquad b_c=P_0b,\qquad P_0u=0.
\]
Here $b_c$ is a surface connection with constant curve coefficients.
Before imposing flatness, it can differ from $b(o)$.

\begin{proposition}\label{mht:prop:exact-slice-action}
The physical action is exactly
\begin{align}
 S_E(b+\eta_0J)
  &=\kappa\int_U\operatorname{tr}(JF_{b_c})
     -\frac12\int_U\tau_E\operatorname{tr}(u\,d_Ju)\notag\\
  &\hspace{1em}
     -\frac{\kappa}{2}\int_{\partial U}\operatorname{tr}(b_cJ).
       \label{mht:eq:exact-slice-action}
\end{align}
The boundary term is absent on a closed surface or with compact support
in its interior. The differential $d_J$ acts on curve coefficients.
In oriented surface coordinates $(s,t)$, the middle term is
$-\int_U\tau_E\operatorname{tr}(u_s d_Ju_t)\,ds\,dt$.

The infinitesimal gauge derivative of the condition $ha=0$, on gauge
parameters $\xi\in V$, is
\begin{equation}\label{mht:eq:gauge-slice-derivative}
 M_J=hd_J|_V=1+h\eta_0\operatorname{ad}_J|_V,\qquad
 M_J^{-1}=h_J\bar\partial|_V.
\end{equation}
\end{proposition}

\begin{proof}
Expand the kinetic and cubic terms in $b_c+u+\eta_0J$.
Every term linear in $u$ vanishes because
$\tau_E(u\eta_0)=\kappa P_0u=0$,
$\tau_E(\bar\partial u)=0$, and $P_0$ commutes with $d_U$.
The terms independent of $u$ are the BF density and its boundary
term from Proposition~\ref{mht:prop:bf-density}.

The remaining cubic coefficient is
$\tau_E\operatorname{tr}(\eta_0Ju^2)$.
Cyclicity identifies it with the
$\eta_0\operatorname{ad}_J$ part of the displayed quadratic form.
The remaining kinetic coefficient is its $\bar\partial$ part.
A product of three surface fields has no curve one-form and
contributes zero. This proves the exact identity for
surface-dependent $J$ as well.

At the slice, the curve component of a gauge variation is $d_J\xi$.
Applying $h$ gives $M_J$. Since $d_JV\subset W$ and
$\bar\partial h=1$ on $W$, one has $\bar\partial M_J=d_J$ on $V$.
Composition of the two isomorphisms $d_J,\bar\partial:V\to W$ gives
the inverse in \eqref{mht:eq:gauge-slice-derivative}.
\end{proof}

The eliminated surface fields therefore have a quadratic action whose
kernel depends on $J$. A quantum elimination must use this same
background dependence both for the gauge determinant and for observable
contractions. The classical splitting alone assigns neither quantity.

\section{The polar first-jet boundary}

The condition $j(f)=0$ makes the polar trace closed. To allow the first
jet to vary, retain it as a boundary variable. Write
\[
 \mathcal A=\Omega^{0,\bullet}(C),\qquad
 B=\mathbb C[\xi]/(\xi^2),\qquad |\xi|=0,\qquad d_B=0.
\]
Here $\xi$ is even, unlike the odd residual class
$\varepsilon\in A_\kappa$. Define differential graded algebra maps
\begin{equation}\label{mht:eq:jet-boundary-maps}
 \begin{split}
 r:\mathcal A\longrightarrow B,&\qquad
 r(f)=f(o)+\xi j(f),\quad r(\eta)=0,\quad |\eta|=1,\\
 \ell:\mathbb C\longrightarrow B,&\qquad \ell(c)=c.
 \end{split}
\end{equation}

The product rule for $j$ proves multiplicativity of $r$. The map
annihilates positive degree, so it commutes with the differentials.
The Neumann algebra is the strict pullback
\begin{equation}\label{mht:eq:neumann-pullback}
 E=\mathcal A\times_B\mathbb C.
\end{equation}
Indeed, the equation $r(f)=\ell(c)$ says $j(f)=0$ and $c=f(o)$.
There is no restriction in degree one.

Let $\tau_{\mathcal A}$ be the polar functional on $\mathcal A^1$,
extended by zero on $\mathcal A^0$, and define
\[
 \tau_B(c_0+\xi c_1)=\kappa c_1.
\]
The polar Stokes identity becomes the relative trace identity
\begin{equation}\label{mht:eq:relative-trace}
 \tau_{\mathcal A}\bar\partial=\tau_Br,\qquad
 \tau_B\ell=0.
\end{equation}
The two traces have degrees $-1$ and $0$, respectively.

\begin{proposition}\label{mht:prop:jet-boundary-pairing}
The pairing $(x,y)\mapsto\tau_B(xy)$ on $B$ is nondegenerate, and
$\ell(\mathbb C)$ is Lagrangian. The same holds after tensoring with
$\mathfrak{gl}_N$ and the matrix trace: the constant matrices form a
Lagrangian Lie subalgebra of $\mathfrak{gl}_N(B)$.
\end{proposition}

\begin{proof}
In the basis $(1,\xi)$ the pairing has matrix
\[
 \kappa\begin{pmatrix}0&1\\1&0\end{pmatrix}.
\]
The constant line is isotropic and equals its orthogonal complement.
For matrices the pairing is
\[
 \kappa\operatorname{tr}(X_0Y_1+X_1Y_0).
\]
Nondegeneracy of the ordinary matrix trace gives the assertion, and
cyclicity makes the pairing invariant under the Lie bracket.
\end{proof}

Thus the Neumann condition selects the constant Lagrangian subalgebra
of the first-jet boundary. This boundary records the integration term
at the double pole without choosing a real boundary circle.

\subsection{The boundary action}

For this subsection assume that $U$ is closed and oriented. Tensor the
relative diagram \eqref{mht:eq:jet-boundary-maps} with
$\Omega^\bullet(U)\otimes\mathfrak g$, retaining the full graded
fields. Let $\mathbf A$ be the total-degree-one bulk superfield,
$D=d_U+\bar\partial$, and $a=r(\mathbf A)$. Put
\[
 T_{\mathcal A}=\int_U\tau_{\mathcal A}\operatorname{tr},\qquad
 T_B=\int_U\tau_B\operatorname{tr},\qquad
 S_{\mathcal A}(\mathbf A)
 =T_{\mathcal A}\left(\frac12\mathbf A D\mathbf A
                                      +\frac13\mathbf A^3\right).
\]

The traces $T_{\mathcal A}$ and $T_B$ have degrees $-3$ and $-2$.
After shifting the field complexes by one, their induced pairings have
degrees $-1$ and $0$. These are the bulk Batalin--Vilkovisky pairing
and the boundary Batalin--Fradkin--Vilkovisky pairing, respectively.

\begin{proposition}\label{mht:prop:polar-boundary-action}
In the cyclic suspension convention for the full fields,
\begin{equation}\label{mht:eq:polar-action-variation}
 \delta S_{\mathcal A}
 =T_{\mathcal A}\bigl(\delta\mathbf A
                     (D\mathbf A+\mathbf A^2)\bigr)
                         -\frac12T_B(a\,\delta a).
\end{equation}
The boundary one-form and Hamiltonian are
\begin{align}
 \Theta_B&=\frac12T_B(a\,\delta a),&
 \omega_B&=\delta\Theta_B,\label{mht:eq:polar-boundary-form}\\
 S_B(a)&=T_B\left(\frac12a\,d_Ua+\frac13a^3\right).
       \label{mht:eq:polar-boundary-hamiltonian}
\end{align}
Writing $a=a_0+\xi a_1$, one obtains
\begin{equation}\label{mht:eq:polar-bfv-bf}
 S_B=\kappa\int_U\operatorname{tr}
                            \bigl(a_1(d_Ua_0+a_0^2)\bigr).
\end{equation}
The boundary condition $a_1=0$ is Lagrangian and preserved by the
cohomological vector field. Its restriction recovers the Neumann
action.
\end{proposition}

\begin{proof}
The kinetic variation is
$\tfrac12T_{\mathcal A}(\delta\mathbf A D\mathbf A
                      +\mathbf A D\delta\mathbf A)$.
Use the identity
\[
 D(\mathbf A\delta\mathbf A)
       =(D\mathbf A)\delta\mathbf A-\mathbf A D\delta\mathbf A
\]
and \eqref{mht:eq:relative-trace}. The resulting endpoint term is
$-\tfrac12T_B(a\delta a)$, and the interior kinetic terms combine.
Graded cyclicity makes the cubic variation
$T_{\mathcal A}(\delta\mathbf A\,\mathbf A^2)$.
This proves \eqref{mht:eq:polar-action-variation} and gives the relative
Hamiltonian identity
$\iota_Q\omega_{\mathcal A}=\delta S_{\mathcal A}+r^*\Theta_B$.

On the boundary, the vector field is $Qa=d_Ua+a^2$.
Associativity and the derivation identity give $Q^2=0$ in this
suspension convention. To expand its Hamiltonian, retain the coefficient
of $\xi$:
\[
 \frac{\kappa}{2}\int_U\operatorname{tr}
           (a_0d_Ua_1+a_1d_Ua_0)
 +\frac{\kappa}{3}\int_U\operatorname{tr}
           (a_0^2a_1+a_0a_1a_0+a_1a_0^2).
\]
Surface integration by parts identifies the two kinetic terms, and
graded cyclicity identifies the three cubic terms. This yields
\eqref{mht:eq:polar-bfv-bf}. Moreover,
\[
 Qa_0=d_Ua_0+a_0^2,\qquad
 Qa_1=d_Ua_1+a_0a_1+a_1a_0.
\]

Hence $a_1=0$ is invariant. Its Lagrangian property follows from
Proposition~\ref{mht:prop:jet-boundary-pairing}. Both $S_B$ and the
boundary one-form restrict to zero there.
\end{proof}

The degrees in this calculation matter. The variables $a_0$ and $a_1$
include ghosts and antifields as well as ordinary surface forms.
Restricting them both to one-forms would erase the boundary Hamiltonian
by surface degree. If $U$ has a boundary, its Stokes term must also be
retained. The polar calculation alone does not give the corner data.

\section{A closed path model for the relative trace}

The pullback \eqref{mht:eq:neumann-pullback} fixes the boundary jet to
be constant. A path from the unrestricted jet to a constant instead
retains the boundary comparison itself. Let $B[t,dt]$ be the polynomial
de Rham algebra of the interval, with $|t|=0$, $|dt|=1$ and
$d_t t=dt$. Define
\begin{equation}\label{mht:eq:relative-path-algebra}
 \widetilde E=\bigl\{(a,b,c)\in
     \mathcal A\oplus B[t,dt]\oplus\mathbb C:
                  b(0)=r(a),\ b(1)=\ell(c)\bigr\}.
\end{equation}
Endpoint evaluation sets $dt=0$. The product is componentwise and the
differential is $D=(\bar\partial,d_t,0)$.

Write a degree-one element as
$y=(\eta,dt(\beta_0(t)+\xi\beta_1(t)),0)$, and set
\begin{equation}\label{mht:eq:relative-path-trace}
 \widetilde\tau(y)=\tau_{\mathcal A}(\eta)
                            +\kappa\int_0^1\beta_1(t)\,dt.
\end{equation}
Extend this functional by zero in degree zero. It is a graded cyclic
trace of degree $-1$. For $x=(a,b,c)$ in degree zero,
\[
 \widetilde\tau(Dx)
 =\tau_Br(a)+\tau_B\bigl(b(1)-b(0)\bigr)
 =\tau_B\ell(c)=0.
\]
Thus the path integral cancels the failure of the bulk trace to be
closed.

\begin{proposition}\label{mht:prop:relative-path-retraction}
The multiplicative inclusion $\iota:E\to\widetilde E$, given by
\[
 \iota(f)=(f,f(o),f(o)),\qquad \iota(\eta)=(\eta,0,0),
\]
is a trace-preserving quasi-isomorphism. It has the following cochain
retraction. For
$x=(a,b_0(t)+\xi b_1(t),c)$ and $y$ as above, put
\begin{align}
 p_0x&=a-f_0j(a),&
 p_1y&=\eta+\eta_0I_1,&
 I_i&=\int_0^1\beta_i(t)\,dt.\label{mht:eq:relative-path-projection}
\end{align}
Let $K$ vanish in degree zero and define
\begin{equation}\label{mht:eq:relative-path-homotopy}
 Ky=\left(-f_0I_1,
  \int_0^t\beta_0(u)\,du
    +\xi\left[-I_1+\int_0^t\beta_1(u)\,du\right],I_0\right).
\end{equation}
Then
\[
 DK+KD=1-\iota p,\qquad p\iota=1,\qquad
 K^2=K\iota=pK=0,\qquad \widetilde\tau=\tau_Ep.
\]
\end{proposition}

\begin{proof}
The zero endpoint of the middle component of $Ky$ is $-\xi I_1$,
which is $r(-f_0I_1)$. The endpoint at one is $I_0$. Thus $Ky$ belongs
to \eqref{mht:eq:relative-path-algebra}.

On degree one, $DKy$ has first component $-\eta_0I_1$ and path
component $dt(\beta_0+\xi\beta_1)$. These are the components of
$y-\iota p_1y$. On degree zero, the endpoint conditions give
\[
 \int_0^1b'_1(t)\,dt=-j(a),\qquad
 \int_0^1b'_0(t)\,dt=c-a(o).
\]
Substitution in \eqref{mht:eq:relative-path-homotopy} gives
$KDx=x-\iota p_0x$. The chain-map identity for $p$ also follows
directly:
\[
 p_1Dx=\bar\partial a-\eta_0j(a)=\bar\partial p_0x.
\]

The identities $p\iota=1$ and $K\iota=0$ follow from the constant
path in the inclusion. Degree gives $K^2=0$, while
$p_0(-f_0I_1)=0$ gives $pK=0$.
Finally, $\tau_E(\eta_0)=\kappa$ proves the trace identity.
\end{proof}

\begin{example}\label{mht:ex:path-adjoint-defect}
This retraction does not preserve the cyclic pairings. For
\[
 v=\frac{|z|^2}{1+|z|^2}\in E^0,
 \qquad y=(0,\xi\,dt,0),
\]
one has
\[
 \widetilde\tau(\iota(v)y)=\kappa v(o)=0,\qquad
 \tau_E(vp_1y)=\tau_E(v\eta_0)=\frac\kappa2.
\]
The last value is the integral
$\kappa\int_0^\infty r(1+r)^{-3}\,dr$.
\end{example}

In fact no projection into unchanged smooth $E$ can be the cyclic
adjoint of $\iota$ on this $y$. Such a projection would give a smooth
form $\eta$ satisfying $\tau_E(v\eta)=\kappa v(o)$ for every
$v\in E^0$. Functions supported away from $o$ belong to $E^0$ and
would force $\eta=0$ away from $o$, hence everywhere by smoothness.
The test $v=1$ then contradicts the required value $\kappa$.

The closed trace on $\widetilde E$ nevertheless gives a classical
cyclic action after tensoring with the surface and matrix algebras.
Its restriction through $\iota$ is the Neumann action. The cochain
retraction proves equivalence of the complexes, but supplies no smooth
cyclic inverse pairing or elliptic propagator on the path algebra.

\section{The residual class is global on the curve}

The global normal form depends on solving a problem over the whole
sphere. A local gauge calculation explains why pole jets alone cannot
replace that step. In this calculation, $w$ is a formal
antiholomorphic variable, independent of $z$. Write a matrix connection
germ as
\[
 a=A(z,w)\,dw,\qquad A(z,w)=\sum_{n\geq0}A_n(z)w^n,
 \qquad A_n(z)\in\operatorname{Mat}_N(\mathbb C[[z]]).
\]

\begin{proposition}\label{mht:prop:local-curve-frame}
There is a unique formal frame $g$ with $g(z,0)=I_N$ and
$\partial_wg=-Ag$. It is an admissible Neumann gauge and sends $a$
to zero. Consequently, every gauge-invariant formal function of this
connection germ alone is constant.
\end{proposition}

\begin{proof}
For $g=\sum_{n\geq0}g_n(z)w^n$, the equation is equivalent to
\begin{equation}\label{mht:eq:local-frame-recursion}
 g_0=I_N,\qquad
 g_{n+1}=-\frac1{n+1}\sum_{j=0}^n A_jg_{n-j}.
\end{equation}
The recursion constructs each coefficient uniquely without commuting
matrix factors. The constant term makes $g$ invertible, and
$g(z,0)=I_N$ gives $j(g)=0$. Its gauge action sends $a$ to zero by
the defining differential equation. An invariant function therefore
takes on every germ its value at zero.
\end{proof}

The same distinction appears linearly. The exact sequence of global
complexes
\[
 0\longrightarrow E\longrightarrow\mathcal A
          \xrightarrow{j}\mathbb C[0]\longrightarrow0
\]
gives $H^1(E)=\mathbb C[\eta_0]$, because a global holomorphic function
has zero derivative and $H^1(\mathcal A)=0$. On a disk, the holomorphic
function $z$ has derivative one, so the analogous local first cohomology
vanishes.

There is an explicit representative invisible to every jet at the
pole. Choose a smooth cutoff $\chi(r)$, equal to one near $r=0$ and
zero outside a small coordinate disk, and define
\begin{equation}\label{mht:eq:annular-residual}
 f_*=z\chi(|z|^2),\qquad
 \eta_*=\bar\partial f_*=z^2\chi'(|z|^2)\,d\bar z.
\end{equation}
Extend $f_*$ by zero to the sphere. The form $\eta_*$ is supported in
an annulus, so all of its jets at $o$ vanish. Nevertheless,
\[
 \tau_E(\eta_*)=-\kappa\int_0^\infty\chi'(r)\,dr=\kappa.
\]

For a constant matrix $J$, the connection $a=\eta_*J$ has normalized
global frame $H=e^{-f_*J}$ and residual matrix $\mathcal J(a)=J$.
This formula holds formally and, for this family, at every finite $J$.
Thus even all local jets of $a$ fail to determine the global residual.
The assertion concerns the curve connection alone. Surface curvature
and additional fields may carry their own local invariants.

\section{Surface-dependent gauge transformations}

Let $D=d_U+\bar\partial$ and let $g$ be a smooth formal Neumann gauge
near the identity, defined on $U\times C$. Set
\[
 \theta=g^{-1}Dg,\qquad \theta_R=(Dg)g^{-1},\qquad
 \operatorname{CS}(A)=\operatorname{tr}
                \left(\frac12 A\,DA+\frac13A^3\right).
\]
The one-forms $\theta$ and $\theta_R$ are conjugate, but their placement
in the finite gauge identity is fixed by the convention
$A^g=g^{-1}Ag+\theta$.

\begin{lemma}\label{mht:lem:finite-transgression}
As a cyclic three-form density,
\begin{equation}\label{mht:eq:finite-transgression}
 \operatorname{CS}(A^g)-\operatorname{CS}(A)
   =-\frac12D\operatorname{tr}(\theta_RA)
                       -\frac16\operatorname{tr}(\theta^3).
\end{equation}
\end{lemma}

\begin{proof}
It is convenient to set $k=g^{-1}$ and
$\vartheta=k^{-1}Dk=-\theta_R$.
The transformed connection is $kAk^{-1}-(Dk)k^{-1}$.
Conjugating it and its derivative by $k^{-1}$ gives
$A-\vartheta$ and
$DA+\vartheta A+A\vartheta-\vartheta^2$.
Here $D\vartheta=-\vartheta^2$.

Substitution in the kinetic and cubic terms, with graded cyclicity,
reduces their difference to
\[
 -\frac12\operatorname{tr}
       (\vartheta^2A+\vartheta DA)
       +\frac16\operatorname{tr}(\vartheta^3)
 =\frac12D\operatorname{tr}(\vartheta A)
       +\frac16\operatorname{tr}(\vartheta^3).
\]
Substitute $\vartheta=-\theta_R$ and use conjugacy of
$\theta_R$ and $\theta$.
\end{proof}

Choose a formal Neumann path $g_u$, $0\leq u\leq1$, from $1$ to $g$.
The parameter $u$ is independent of the surface coordinates.
Put
\[
 v_u=g_u^{-1}\partial_ug_u,\qquad \theta_u=g_u^{-1}Dg_u.
\]
Path integrals are interpreted coefficientwise in the formal
amplitude.

\begin{proposition}\label{mht:prop:surface-boundary-functional}
For $S_E(A)=\int_U\tau_E\operatorname{CS}(A)$,
\begin{align}
 S_E(A^g)-S_E(A)&=\mathcal B(A,g),\label{mht:eq:boundary-difference}\\
 \mathcal B(A,g)
 &=-\frac12\int_{\partial U}\tau_E\operatorname{tr}(\theta_RA)
   -\frac12\int_{\partial U}\int_0^1
               \tau_E\operatorname{tr}(v_u\theta_u^2)\,du.
       \label{mht:eq:boundary-functional}
\end{align}
This integrated functional is independent of the chosen formal path
and satisfies
\begin{equation}\label{mht:eq:boundary-cocycle}
 \mathcal B(A,gh)=\mathcal B(A,g)+\mathcal B(A^g,h).
\end{equation}
These assertions concern small formal gauges defined on $U\times C$.
\end{proposition}

\begin{proof}
The Maurer--Cartan identity and its parameter derivative are
\[
 D\theta_u=-\theta_u^2,\qquad
 \partial_u\theta_u=Dv_u+[\theta_u,v_u].
\]
Graded cyclicity removes the commutator in the trace, while
$D(\theta_u^2)=0$. Hence
\[
 \partial_u\operatorname{tr}(\theta_u^3)
             =3D\operatorname{tr}(v_u\theta_u^2).
\]
Integrate along the path and substitute in
\eqref{mht:eq:finite-transgression}.

The surface coefficients of $\theta_u,\theta_R,A$, and the function
$v_u$, have zero first holomorphic jet. Their degree-zero curve products
are therefore in $E^0$. Polar Stokes contributes zero.
Surface Stokes leaves exactly the oriented boundary expression
\eqref{mht:eq:boundary-functional}.

The action difference depends only on the endpoints, proving
independence of the integrated functional from the path.
Finally, $(A^g)^h=A^{gh}$ makes the action differences telescope and
gives \eqref{mht:eq:boundary-cocycle}.
\end{proof}

This construction does not assign periods to large gauge
transformations. On its formal domain, it gives a classical boundary
compensator:
\[
 \widetilde S(A,g)=S_E(A)+\mathcal B(A,g)=S_E(A^g).
\]
Under $(A,g)\mapsto(A^k,k^{-1}g)$, the connection $A^g$ stays fixed.
Thus the group-valued edge field $g$ makes $\widetilde S$ invariant.
Its identification with a geometric closed-string field would require
a separate map.

\subsection{A nonzero cubic boundary response}

Take $U=[0,1]^2$ with coordinates $(s,t)$ and orientation $ds\wedge dt$.
On the sphere set
\[
 f=\frac{z^2}{(1+|z|^2)^2},\qquad
 g_1=\frac{\bar z}{1+|z|^2},\qquad
 v=\frac{|z|^2}{1+|z|^2}.
\]
These are smooth Neumann functions. Let
$E=E_{12}$, $F=E_{21}$, $Z=[E,F]$, and define
\[
 X=sfE+tg_1F+st\,vZ,\qquad g_\lambda=\exp(\lambda X),
\]
where $\lambda$ is a formal amplitude. Since $j(X)=0$,
$g_\lambda$ is an admissible gauge.

Matrix multiplication, with surface forms before $d\bar z$, gives
\begin{align*}
 \tau_E\operatorname{tr}((DX)^3)
 &=6st\,\tau_E\bigl(fg_1\bar\partial v-fv\bar\partial g_1
                              -g_1v\bar\partial f\bigr)\,ds\,dt\\
 &=12st\,\tau_E(fg_1\bar\partial v)\,ds\,dt.
\end{align*}
The second equality uses polar Stokes for $fg_1v$, whose first jet is
zero. The remaining integral is explicit:
\[
 fg_1\bar\partial v
      =\frac{z^2|z|^2}{(1+|z|^2)^5}\,d\bar z,\qquad
 \tau_E(fg_1\bar\partial v)
      =-\kappa\int_0^\infty\frac{r}{(1+r)^5}\,dr
      =-\frac{\kappa}{12}.
\]
Thus $\tau_E\operatorname{tr}((DX)^3)=-\kappa st\,ds\,dt$, and
\begin{equation}\label{mht:eq:nonzero-wess-zumino}
 S_E(0^{g_\lambda})
   =-\frac{\lambda^3}{6}\int_U\tau_E\operatorname{tr}((DX)^3)
             +O(\lambda^4)
   =\frac{\kappa\lambda^3}{24}+O(\lambda^4).
\end{equation}
For the path $\exp(u\lambda X)$, the leading boundary primitive is
$-\lambda^3\tau_E\operatorname{tr}(X(DX)^2)/6$.
Its derivative is the leading bulk density, so the four oriented
boundary edges give the same coefficient.

A frame depending only on the curve has $\theta^2=0$ by curve degree.
Equation~\eqref{mht:eq:nonzero-wess-zumino} shows why that simplification
fails for a surface-dependent frame. The residual comparison must retain
the boundary functional and its composition law. A quantum comparison
must in addition account for the background-dependent quadratic
complement in \eqref{mht:eq:exact-slice-action}.
